{{دیوان
| عنوان = تضمین بر شعرِ انیسی شاملوؔ
| ماخذ = ویکی ماخذ
| ماخذ_لنک = https://ur.wikisource.org/wiki/%D8%AA%D8%B6%D9%85%DB%8C%D9%86_%D8%A8%D8%B1_%D8%B4%D8%B9%D8%B1%D9%90_%D8%A7%D9%86%DB%8C%D8%B3%DB%8C_%D8%B4%D8%A7%D9%85%D9%84%D9%88
| مقابلہ = اقبال سائبر لائبریری (اقبال اکادمی پاکستان)
| مقابلہ_لنک = https://www.iqbalcyberlibrary.net
}}
<poem>
ہمیشہ صورتِ بادِ سحَر آوارہ رہتا ہوں
محبّت میں ہے منزل سے بھی خوشتر جادہ پیمائی

دلِ بے تاب جا پہنچا دیارِ پیرِ سنجر میں
میسّر ہے جہاں درمانِ دردِ ناشکیبائی

ابھی ناآشنائے لب تھا حرفِ آرزو میرا
زباں ہونے کو تھی منّت پذیرِ تابِ گویائی

یہ مرقد سے صدا آئی، حرم کے رہنے والوں کو
شکایت تجھ سے ہے اے تارکِ آئینِ آبائی!

ترا اے قیس کیونکر ہوگیا سوزِ درُوں ٹھنڈا
کہ لیلیٰ میں تو ہیں اب تک وہی اندازِ لیلائی

نہ تخمِ ’لا الہ‘ تیری زمینِ شور سے پھُوٹا
زمانے بھر میں رُسوا ہے تری فطرت کی نازائی

تجھے معلوم ہے غافل کہ تیری زندگی کیا ہے
کُنِشتی ساز، معمورِ نوا ہائے کلیسائی

ہُوئی ہے تربیت آغوشِ بیت اللہ میں تیری
دلِ شوریدہ ہے لیکن صنم خانے کا سودائی

“وفا آموختی از ما، بکارِ دیگراں کر دی
ربودی گوہرے از ما نثارِ دیگراں کر دی”
</poem>
